\section{What Happens in the Operating Room}

\subsection{The room, the cabinet, and the cable that is not connected}

\draftinstr{Roughly 1,200 characters introducing Figure 17. State the three
things that cross the network boundary before a procedure, a signed model build,
a protocol version, and the pinned registry entry; state the two that cross
after it, reports to the FDA and the IRB; and state that nothing crosses during
it. Cite \texttt{FDACyber2026} for the cybersecurity expectation and
\texttt{onpremwhippl} for the on-premises design. Note that the stack topology
is adapted from the author's NSCLC single-patient journey \cite{paipatientjr}
and re-scoped to PDAC: eight arms rather than four, three anastomoses rather
than one bronchial closure, and a vascular exclusion envelope the thoracic
configuration has no equivalent of.}

\begin{pafig}[diagrams (python)-type: deployment, full page]
\node[dgtiled] (a) at (0,0) {};
\glyphserver{a}{protowhite}
\node[dglabel,anchor=north] at ([yshift=-5.4mm]a.center) {Figure 17 slot, OR and on-prem stack};
\node[dglabel,anchor=north] at (0,-1.6) {\ttfamily\tiny Source\\ diagrams-python/\\ fig\_\pfb 17\_\pfb operating\_\pfb room\_\pfb stack.py};
\end{pafig}
\vspace{-0.7cm}
\figcaption{Figure 17. The operating room and the on-premises stack drawn as a deployment
architecture, with the sterile field, the non-sterile side, the cabinet, the evidence\\
store, and the network boundary that carries no intra-procedural route at all.}

\draftinstr{Replace the Figure 17 slot with the full-page Diagrams
(Python)-type deployment from
\texttt{diagrams-\pfb python/\pfb fig\_\pfb 17\_\pfb operating\_\pfb room\_\pfb stack.py}. Five
\texttt{dgcluster} containers in the two-column arrangement given in that file;
every tile a \texttt{\textbackslash dgnode} with the named pictogram; the absent
cable interrupted by a \texttt{\textbackslash pxmark} with its label in a
white-filled node above the cross.}

\subsection{The eight arms, and what limits them}

\draftinstr{Roughly 1,300 characters. Per-arm tip force at most 3 N, cumulative
across arms at most 18 N, positional tolerance 2 mm, force reproducibility
0.5 N, a 10 kHz heartbeat broadcast bus, and a vascular no-fly corridor 4 mm
either side of the superior mesenteric and portal veins. State the worst
observed value from the Phase 0 simulation set beside each cap, and state that a
simulated worst case is not a human worst case. Cite \texttt{pdac060s2030} and
\texttt{iec8060127}.}

\subsection{Who is controlling the operation}

\draftinstr{Roughly 1,400 characters, and treat this as the most important
subsection in the paper. This is the single most frequently raised concern in
every surveyed source \cite{WuSemiAI2026, BrarRAS2024X, Jauniaux2025}. The model
proposes; a deterministic gate tests the proposal against the force caps and the
no-fly corridor; the surgeon approves, modifies, or declines, and declining
requires no justification; the signed command reaches the arms. There is no
message from the model to the arms. Cite \texttt{LeeAuto2024X} for the published
finding that no FDA-cleared surgical robot exceeds task-level autonomy under
supervision, and \texttt{FDARobot2022}.}

\begin{pafig}[mermaid-type: sequenceDiagram]
\node[mmactor] (a) at (0,0) {Figure 19 slot};
\node[mmin] (b) at (0,-2.0) {\ttfamily\tiny Source\\ mermaid/\\ fig-\pfb 19-\pfb advise-\pfb approve-\pfb execute.md};
\draw[mmedge] (a) -- (b);
\end{pafig}
\vspace{-0.7cm}
\figcaption{Figure 19. One operative step message by message across five lifelines, with the
deterministic gate between advice and approval, and no arrow at all from the model\\
to the arms, which is the safety argument this figure makes by an absence.}

\draftinstr{Replace the Figure 19 slot with the mermaid-type sequence from
\texttt{mermaid/\pfb fig-\pfb 19-\pfb advise-\pfb approve-\pfb execute.md}. Five \texttt{mmactor} headers
at 3.6 cm pitch; \texttt{mmlife} lifelines from $-0.6$ to $-12.6$;
\texttt{mmact} activation bars; self-messages as 4 mm right-hand loops at
looseness 6; the \texttt{alt} fragment with a corner tab; the two closing notes
at $y = -11.2$ and $-12.2$ so they do not collide with the last message. Figure
18 precedes this figure in reading order and is placed above it.}

\subsection{How fast stopping actually is}

\draftinstr{Roughly 1,100 characters introducing Figure 18. Three budgets: a
0.1 ms heartbeat period, a 3 ms cross-arm halt, and a 500 ms system-wide
guarantee. Give the human referent: a blink is roughly 100 to 150 ms, so every
arm is held before a blink completes. Name the three parties who may press stop:
the operating surgeon, the second operator station, and the independent safety
monitor. Cite \texttt{ieee7009}.}

\begin{pafig}[plantuml-type: timing, @startuml]
\node[umlkey] (a) at (0,0) {Figure 18 slot\\e-stop timing budget};
\node[umlbox] (b) at (0,-1.9) {\ttfamily\tiny Source\\ plantuml/\\ fig-\pfb 18-\pfb estop-\pfb timing-\pfb budget.puml};
\draw[umlarrow] (a) -- (b);
\end{pafig}
\vspace{-0.7cm}
\figcaption{Figure 18. The emergency-stop chain on a logarithmic millisecond axis, from the
instant the surgeon presses stop through the bus broadcast and arm deceleration to\\
zero tip force, with all three protocol budget lines marked across every timeline.}

\draftinstr{Replace the Figure 18 slot with the PlantUML-type timing diagram
from \texttt{plantuml/\pfb fig-\pfb 18-\pfb estop-\pfb timing-\pfb budget.puml}. Five stacked 14 cm
bands at 1.6 cm pitch; a logarithmic axis from 0.1 to 1000 ms with decade ticks;
state segments as 4 mm rectangles filled \texttt{protoblue} while asserted,
\texttt{pablue2} while armed, \texttt{pagraym} while idle; budget markers as
0.8pt dashed verticals spanning all five bands with the label in a white-filled
node. Place this figure before Figure 19 on the page.}

\subsection{The vascular no-fly corridor and the force envelope}

\draftinstr{Roughly 900 characters introducing Figure 21. A force cap does not
prevent an arm entering a vascular corridor, and a corridor does not cap the
force applied inside permitted tissue. The pairing is the safety property, and
the figure draws them as sibling containers inside one envelope for that reason.}

\begin{pafig}[d2-type: container with measurement grid]
\node[d2key] (a) at (0,0) {Figure 21 slot\\force and no-fly envelope};
\node[d2box] (b) at (0,-1.8) {\ttfamily\tiny Source\\ d2/\\ fig-\pfb 21-\pfb force-\pfb nofly-\pfb envelope.d2};
\draw[d2edge] (a) -- (b);
\end{pafig}
\vspace{-0.7cm}
\figcaption{Figure 21. The operative envelope drawn as two sibling containers, the force
limits with their observed headroom on the left and the vascular no-fly geometry on\\
the right, joined by the deterministic gate that rejects any plan violating either.}

\draftinstr{Replace the Figure 21 slot with the D2-type envelope figure from
\texttt{d2/\pfb fig-\pfb 21-\pfb force-\pfb nofly-\pfb envelope.d2}. One outer \texttt{d2cont}, two
\texttt{d2cont2} sub-containers 1.6 cm apart; the force sub-container as a
5-row measurement grid using \texttt{\textbackslash hbarrow} with the cap as a
dashed vertical and the observed value as a filled bar; the geometry
sub-container as a schematic with two 1.4 mm vessel bands, a hatched
\texttt{pagraym} corridor, and eight arm-tip discs, none inside the corridor;
the invariant strip in \texttt{d2dark} along the foot.}

\subsection{What could fail, and what catches it}

\draftinstr{Roughly 1,200 characters introducing Figure 20. Four intermediate
events must all occur for harm to reach the participant, which is what the AND
gate at the apex means. Name the four barriers, and name the one gap: registration
drift beyond 2 mm is caught by the deterministic gate alone. Do not omit the gap.}

\begin{pafig}[graphviz-type: fault tree, AND and OR gates]
\node[gvboxd] (a) at (0,0) {Figure 20 slot\\hazard and barrier fault tree};
\node[gvbox] (b) at (0,-2.0) {\ttfamily\tiny Source\\ graphviz/\\ fig-\pfb 20-\pfb hazard-\pfb barrier-\pfb fault-\pfb tree.dot};
\draw[gvedge] (a) -- (b);
\end{pafig}
\vspace{-0.7cm}
\figcaption{Figure 20. The hazard fault tree with the top event behind an AND gate, so that
harm requires all four intermediate events together, and with the single-barrier gap\\
at registration drift marked in medium-dark gray rather than quietly omitted.}

\draftinstr{Replace the Figure 20 slot with the graphviz-type fault tree from
\texttt{graphviz/\pfb fig-\pfb 20-\pfb hazard-\pfb barrier-\pfb fault-\pfb tree.dot}. Use
\texttt{\textbackslash umlgateand} and \texttt{\textbackslash umlgateor} for the
five gates; ten basic events as 9 mm \texttt{gvcircle} nodes on the lowest rank;
barrier annotations in the right margin attached by 0.5pt dashed leaders so they
never sit between two events; the two cross-branch edges at looseness 0.9.}

\subsection{Conversion to open surgery}

\draftinstr{Roughly 800 characters. Conversion is a recorded outcome, not a
failure of the study or of the participant, it does not end participation, and
the decision belongs to the operating surgeon alone. Give the comparator
conversion rate context from \texttt{DutchCohort2025} and state that backup
instruments and an open tray are in the room for every procedure.}
